% ======================================================================
% Encoding, Sprache, Schriften
% ======================================================================
\usepackage[utf8]{inputenc} % Bei aktuellem LaTeX oft entbehrlich, aber kompatibel.
\usepackage[T1]{fontenc}    % T1-Font-Encoding für korrekte Trennung/Umlaute.
\usepackage[ngerman,shorthands=off]{babel} % Deutsch + Deaktivierung von babel-shorthands.
\usepackage{csquotes}       % Kontextabhängige Anführungszeichen (empfohlen mit biblatex).

% Serifenschrift (Times-ähnlich) für Fließtext und Mathematik
\usepackage{newtxtext,newtxmath}

% Sans-Serif für Überschriften/Titel (Arial-Ersatz: Arimo)
\usepackage[sfdefault,scaled=0.95]{arimo}
\renewcommand{\familydefault}{\rmdefault} % Fließtext bleibt serif (Times-ähnlich).

% Einheitliche Abstände für Display-Matheumgebungen
\AtBeginDocument{%
  % Formeln als eigener Absatz (typisch mit Leerzeile davor/danach)
  \setlength{\abovedisplayskip}{10pt plus 4pt minus 5pt}%
  \setlength{\belowdisplayskip}{10pt plus 4pt minus 5pt}%
  % Formeln „kurz“ nach Einleitung im Fließtext (z. B. „… wie folgt:“)
  \setlength{\abovedisplayshortskip}{6pt plus 3pt minus 3pt}%
  \setlength{\belowdisplayshortskip}{6pt plus 3pt minus 3pt}%
  % Zeilenabstand innerhalb von align-Umgebungen
  \setlength{\jot}{2pt}%
}

% ======================================================================
% Überschriften (titlesec): Kapitel-/Abschnittsdesign
% ======================================================================
\usepackage{titlesec}

% Kapitel: 14 pt, Sans-Serif fett
\titleformat{\chapter}[hang]
  {\sffamily\bfseries\fontsize{14pt}{16pt}\selectfont}
  {\thechapter}{1em}{}
\titleformat{name=\chapter,numberless}[hang]
  {\sffamily\bfseries\fontsize{14pt}{16pt}\selectfont}{}{0pt}{}
\titlespacing*{name=\chapter,numberless}{0pt}{2pt}{10pt}

% Section/Subsection: 12 pt, Sans-Serif fett
\titleformat{\section}[hang]
  {\sffamily\bfseries\fontsize{12pt}{14pt}\selectfont}
  {\thesection}{1em}{}
\titleformat{\subsection}[hang]
  {\sffamily\bfseries\fontsize{12pt}{14pt}\selectfont}
  {\thesubsection}{1em}{}

% Subsubsection: 12 pt, Sans-Serif normal
\titleformat{\subsubsection}[hang]
  {\sffamily\fontsize{12pt}{14pt}\selectfont}
  {\thesubsubsection}{1em}{}

% ======================================================================
% Kopf- und Fußzeilen (KOMA-Script)
% ======================================================================
\usepackage[
  automark,          % Automatische Marken (Kapitel/Abschnitt) setzen
  headsepline,       % Linie unter der Kopfzeile
  plainheadsepline,  % Linie auch auf plain-Seiten
  ilines,            % Innenliegende Linienführung
  headwidth=text     % Kopfzeile auf Textbreite
]{scrlayer-scrpage}

% ======================================================================
% Grafiken
% ======================================================================
\usepackage{graphicx}
\usepackage{adjustbox}
\usepackage{wrapfig}
\graphicspath{{Bilder/}} % Standardsuchpfad für Grafiken

% ======================================================================
% Mathematik & Symbole
% ======================================================================
\usepackage{amsmath,amsfonts}
\let\openbox\relax % resolve clash: newtxmath already defines \openbox
\usepackage{amsthm}
\usepackage{dsfont}
\usepackage{nicefrac}

% ======================================================================
% Index & Nomenklatur (Abkürzungsverzeichnis)
% ======================================================================
\usepackage{makeidx}
\usepackage[intoc]{nomencl} % Nomenklatur ins Inhaltsverzeichnis aufnehmen
\let\abbrev\nomenclature
\renewcommand{\nomname}{Abkürzungsverzeichnis}
\setlength{\nomlabelwidth}{.25\hsize}
\renewcommand{\nomlabel}[1]{#1 \dotfill}
\setlength{\nomitemsep}{-\parsep}

% ======================================================================
% Tabellen: Zeilenabstände
% ======================================================================
\usepackage{setspace}
\AtBeginEnvironment{tabular}{\singlespacing} % Tabellen stets einzeilig setzen
\renewcommand{\arraystretch}{1.0}            % Zeilenhöhe innerhalb von Tabellen

% ======================================================================
% Sonstige Pakete
% ======================================================================
\usepackage{rotating}
\usepackage{float} % [H] nur sparsam verwenden (Layout kann „festkleben“).

% ======================================================================
% Quellcode (listings)
% ======================================================================
\usepackage{listings}
\usepackage{xcolor}

% Farbdefinitionen für Listings-Stil
\definecolor{codegreen}{rgb}{0,0.6,0}
\definecolor{codegray}{rgb}{0.5,0.5,0.5}
\definecolor{codepurple}{rgb}{0.58,0,0.82}
\definecolor{backcolour}{rgb}{0.95,0.95,0.92}

% Standardstil für Quellcode-Blöcke
\lstdefinestyle{mystyle}{
  backgroundcolor=\color{backcolour},
  commentstyle=\color{codegreen}\itshape,
  keywordstyle=\bfseries,
  numberstyle=\tiny\color{codegray},
  stringstyle=\ttfamily,
  basicstyle=\ttfamily\footnotesize,
  breakatwhitespace=false,
  breaklines=true,
  captionpos=b,
  keepspaces=true,
  numbers=left,
  numbersep=5pt,
  showspaces=false,
  showstringspaces=false,
  showtabs=false,
  tabsize=2
}
\lstset{style=mystyle}

% Monospace-Schrift (Courier) für Listings
\usepackage{courier}

% Wörtliche Zitate: kursiv (für längere Zitatblöcke)
\newenvironment{quoteitalic}{\begin{quote}\itshape}{\end{quote}}

% ======================================================================
% Verzeichnisse (tocloft): Abbildungen, Tabellen, Formeln, Anhang
% ======================================================================
\usepackage{tocloft}

% Einzüge in Abbildungs- und Tabellenverzeichnis entfernen
\setlength{\cftfigindent}{0pt}
\setlength{\cfttabindent}{0pt}

% Titelabstände der Verzeichnisse (wie bei nummerlosen Kapiteln: 2pt/10pt)
\setlength{\cftbeforetoctitleskip}{2pt}
\setlength{\cftaftertoctitleskip}{10pt}
\setlength{\cftbeforeloftitleskip}{2pt}
\setlength{\cftafterloftitleskip}{10pt}
\setlength{\cftbeforelottitleskip}{2pt}
\setlength{\cftafterlottitleskip}{10pt}

% Eigenes Formelverzeichnis
\newlistof{formeln}{loe}{Formelverzeichnis}
\setlength{\cftbeforeloetitleskip}{2pt}
\setlength{\cftafterloetitleskip}{10pt}
\renewcommand{\cftformelnpresnum}{}
\renewcommand{\cftformelnaftersnum}{\enspace}
\setlength{\cftformelnnumwidth}{2.25em}

% Hilfsmakro: Formel in das Formelverzeichnis eintragen
\newcommand{\eqcaption}[1]{%
  \addcontentsline{loe}{formeln}{\protect\numberline{\theequation}#1}%
}

% Überschriftenschrift der Verzeichnisse (analog Kapitelstil)
\renewcommand{\cfttoctitlefont}{\sffamily\bfseries\fontsize{14pt}{16pt}\selectfont}
\renewcommand{\cftloftitlefont}{\sffamily\bfseries\fontsize{14pt}{16pt}\selectfont}
\renewcommand{\cftlottitlefont}{\sffamily\bfseries\fontsize{14pt}{16pt}\selectfont}
\renewcommand{\cftloetitlefont}{\sffamily\bfseries\fontsize{14pt}{16pt}\selectfont}

% Anhangsverzeichnis
\newlistof{anhang}{loah}{Anhangsverzeichnis}
\setlength{\cftbeforeloahtitleskip}{2pt}
\setlength{\cftafterloahtitleskip}{10pt}
\renewcommand{\cftanhangpresnum}{Anhang~}
\renewcommand{\cftanhangaftersnum}{\enspace}
\setlength{\cftanhangnumwidth}{7em}
\renewcommand{\cftloahtitlefont}{\sffamily\bfseries\fontsize{14pt}{16pt}\selectfont}

% ======================================================================
% Hyperref & PDF-Metadaten
% ======================================================================
\usepackage[
  bookmarks,
  bookmarksopen=true,
  pdftitle={\titel},
  pdfauthor={\autor},
  pdfcreator={\autor},
  pdfsubject={\titel},
  pdfkeywords={\titel},
  colorlinks=true,
  linkcolor=black, anchorcolor=black, citecolor=black,
  filecolor=black, menucolor=black, urlcolor=black,
  plainpages=false, pdfpagelabels, hypertexnames=false, linktocpage
]{hyperref}
\usepackage{url}

% ======================================================================
% Bibliographie (biblatex, APA)
% ======================================================================
\usepackage[
  backend=biber,
  style=apa,
  sorting=nyt,
  natbib=false,
  maxcitenames=2,
  maxbibnames=20
]{biblatex}
\DeclareLanguageMapping{ngerman}{ngerman-apa}
\addbibresource{quellen.bib}

% APA-Layout: Hängender Einzug und Abstand zwischen Einträgen
\setlength{\bibhang}{1.27cm}
\setlength{\bibitemsep}{\baselineskip}

% ======================================================================
% Captions & Float-Abstände
% ======================================================================
\usepackage[font={bf,small},labelfont=bf,justification=centering]{caption}

% Abstände zwischen Objekt und Beschriftung (moderat kompakt)
\captionsetup{skip=6pt, belowskip=0pt}
\captionsetup[table]{skip=4pt}

% Abstände rund um Floats (kompakt, aber layoutstabil)
\setlength{\textfloatsep}{10pt plus 2pt minus 4pt} % Float oben/unten auf Seite
\setlength{\intextsep}{8pt plus 2pt minus 3pt}     % Float im Textfluss ([h]/[H])
\setlength{\floatsep}{8pt plus 2pt minus 3pt}      % Zwischen zwei Floats

\DeclareCaptionLabelFormat{anhangnonum}{#1}

% ======================================================================
% Tabellen & Layout-Helfer
% ======================================================================
\usepackage{longtable}
\usepackage{array}
\usepackage{ragged2e}
\usepackage{lscape}
\usepackage{booktabs}
\usepackage{supertabular}
\usepackage[table,xcdraw]{xcolor}
\newcolumntype{w}[1]{>{\raggedleft\arraybackslash}p{#1}}

% Listen und Fußnoten in Tabellen
\usepackage{paralist}
\usepackage{tablefootnote}

% ======================================================================
% Algorithmen
% ======================================================================
\usepackage[linesnumbered,ruled]{algorithm2e}
\SetKwInput{KwInput}{Input}
\SetKwInput{KwOutput}{Output}

% ======================================================================
% Hinweis zur Layoutstabilität
% ----------------------------------------------------------------------
% Keine Neu-Definition von figure/table/equation-Umgebungen, da dies
% Abstände und Nummerierungen unerwartet beeinflussen kann.
% ======================================================================
